\chapter{Martingales and Optimal Stopping}
\label{ch:mart}

A martingale is the mathematical model of a fair game: whatever has happened so far,
the expected future fortune equals the present fortune. The concept is simple, but it is
one of the most powerful tools in probability. It gives short solutions to hitting-time
problems (gambler's ruin, waiting for patterns), sharp concentration inequalities, almost
sure convergence theorems, and the theory of optimal stopping.

Throughout, $X_1,X_2,\dots$ are random variables describing the randomness revealed
step by step, and $\F_n$ denotes ``the information in $X_1,\dots,X_n$'' (with $\F_0$ the
trivial information). We write $\Exp[Y\given\F_n]$ for $\Exp[Y\given X_1,\dots,X_n]$, the
conditional expectation given the vector $(X_1,\dots,X_n)$, defined as in
Chapter~\ref{ch:condexp}. It satisfies the same properties: linearity, taking out what is
known, independence, the tower property
$\Exp[\Exp[Y\given\F_n]\given\F_m]=\Exp[Y\given\F_m]$ for $m\le n$, and the defining
property $\Exp[g(X_1,\dots,X_n)Y]=\Exp[g(X_1,\dots,X_n)\Exp[Y\given\F_n]]$.

\section{Martingales}

\begin{definition}[Martingale]
A sequence $(M_n)_{n\ge0}$ is a \emph{martingale} if
\begin{enumerate}[label=(\roman*)]
\item $M_n$ is determined by $X_1,\dots,X_n$ (\emph{adapted}),
\item $\Exp\abs{M_n}<\infty$, and
\item $\Exp[M_{n+1}\given\F_n]=M_n$ for all $n$.
\end{enumerate}
It is a \emph{submartingale} if $\Exp[M_{n+1}\given\F_n]\ge M_n$ (a favourable game) and a
\emph{supermartingale} if $\le$ (an unfavourable game).
\end{definition}

Taking expectations (tower property), $\Exp[M_n]=\Exp[M_0]$ for a martingale, and by the
tower property $\Exp[M_n\given\F_m]=M_m$ for all $m\le n$. The \emph{increments}
$D_n=M_n-M_{n-1}$ satisfy $\Exp[D_n\given\F_{n-1}]=0$.

\begin{example}[Basic martingales]\label{ex:martingales}
\begin{enumerate}[label=(\arabic*)]
\item \emph{Random walk.} If the $X_i$ are independent with mean $\mu$ and $S_n=\sum_{i\le n}X_i$,
then $S_n-n\mu$ is a martingale: $\Exp[S_{n+1}-(n+1)\mu\given\F_n]=S_n+\mu-(n+1)\mu$.
\item \emph{Squared walk.} If $\Exp[X_i]=0$ and $\Var(X_i)=\sigma^2$, then $S_n^2-n\sigma^2$
is a martingale: $\Exp[S_{n+1}^2\given\F_n]=S_n^2+2S_n\Exp[X_{n+1}]+\sigma^2$.
\item \emph{Exponential martingale.} If $M(\theta)=\Exp[e^{\theta X_i}]<\infty$, then
$e^{\theta S_n}/M(\theta)^n$ is a martingale. For the biased simple walk
($\Prob(X=1)=p$), choosing $\theta=\ln(q/p)$ gives $M(\theta)=p\cdot\frac qp+q\cdot\frac pq=1$,
so $(q/p)^{S_n}$ is a martingale.
\item \emph{Products.} If $Y_i\ge0$ are independent with $\Exp[Y_i]=1$, then
$\prod_{i\le n}Y_i$ is a martingale.
\item \emph{Doob martingale.} For any integrable $Y$, $M_n=\Exp[Y\given\F_n]$ is a martingale
(tower property). Learning information step by step, our estimate of $Y$ is a martingale.
\item \emph{P\'olya's urn.} An urn has $r_0$ red and $b_0$ blue balls; at each step a ball is
drawn, returned, and one extra ball of the same colour added. If $R_n$, $B_n$ are the
numbers after $n$ steps, then the fraction $F_n=R_n/(R_n+B_n)$ is a martingale. Indeed
$R_n+B_n=r_0+b_0+n=:T_n$ is deterministic, and
\[
\Exp[R_{n+1}\given\F_n]=R_n+\frac{R_n}{T_n}=R_n\frac{T_n+1}{T_n},\quad\text{so}\quad
\Exp[F_{n+1}\given\F_n]=\frac{R_n(T_n+1)/T_n}{T_n+1}=F_n.
\]
\item \emph{Branching process.} If $Z_n$ is the generation size with mean offspring $m$,
then $Z_n/m^n$ is a martingale.
\end{enumerate}
\end{example}

\subsection*{You cannot beat a fair game}

\begin{definition}[Predictable process]
A process $(H_n)_{n\ge1}$ is \emph{predictable} if $H_n$ is determined by
$X_1,\dots,X_{n-1}$: the bet at time $n$ may use only the past.
\end{definition}

\begin{theorem}[Martingale transform]\label{thm:transform}
If $M$ is a martingale and $H$ is predictable and bounded, then
$(H\cdot M)_n=\sum_{k=1}^nH_k(M_k-M_{k-1})$ is a martingale, with $\Exp[(H\cdot M)_n]=0$.
\end{theorem}
\begin{proof}
$H_k$ is $\F_{k-1}$-measurable and bounded, so by ``taking out what is known'',
$\Exp[H_k(M_k-M_{k-1})\given\F_{k-1}]=H_k\Exp[M_k-M_{k-1}\given\F_{k-1}]=0$. Hence
$\Exp[(H\cdot M)_{n+1}\given\F_n]=(H\cdot M)_n$.
\end{proof}

Interpreted as betting: $H_k$ is the stake on round $k$ (decided in advance), $M_k-M_{k-1}$
is the gain per unit staked, and the total winnings are still a fair game. No predictable
betting system can create an advantage.

\section{Stopping times and the optional stopping theorem}

\begin{definition}[Stopping time]
A random time $\tau\in\{0,1,2,\dots\}\cup\{\infty\}$ is a \emph{stopping time} if for every
$n$ the event $\{\tau=n\}$ is determined by $X_1,\dots,X_n$: the decision to stop at time $n$
uses only the information available at time $n$.
\end{definition}

For example, ``the first time the walk hits $10$'' is a stopping time; ``the last time before
$100$ that the walk is at its maximum'' is not (it needs the future).

\begin{theorem}[Optional stopping, bounded version]\label{thm:ost0}
If $M$ is a martingale and $\tau$ a stopping time with $\tau\le K$ for a constant $K$, then
$\Exp[M_\tau]=\Exp[M_0]$.
\end{theorem}
\begin{proof}
Let $H_k=\indic_{\{\tau\ge k\}}=1-\indic_{\{\tau\le k-1\}}$, which is determined by
$X_1,\dots,X_{k-1}$, so $H$ is predictable and bounded. Then
$M_{\tau\wedge n}=M_0+\sum_{k=1}^nH_k(M_k-M_{k-1})=M_0+(H\cdot M)_n$, a martingale by
Theorem~\ref{thm:transform}, so $\Exp[M_{\tau\wedge n}]=\Exp[M_0]$ for all $n$. Take
$n=K$.
\end{proof}

\begin{theorem}[Optional stopping, general version]\label{thm:ost}
Let $M$ be a martingale and $\tau$ a stopping time with $\tau<\infty$ almost surely.
Then $\Exp[M_\tau]=\Exp[M_0]$ if any of the following holds:
\begin{enumerate}[label=(\alph*)]
\item $\tau$ is bounded;
\item $\abs{M_{\tau\wedge n}}\le K$ for all $n$ (the stopped process is bounded);
\item $\Exp[\tau]<\infty$ and the increments are bounded: $\abs{M_n-M_{n-1}}\le c$.
\end{enumerate}
\end{theorem}
\begin{proof}
(a) is Theorem~\ref{thm:ost0}. In (b) and (c) we have $\Exp[M_{\tau\wedge n}]=\Exp[M_0]$ and
$M_{\tau\wedge n}\to M_\tau$ almost surely (as $\tau<\infty$). In (b) the sequence is bounded,
so the bounded convergence theorem gives $\Exp[M_{\tau\wedge n}]\to\Exp[M_\tau]$. In (c),
$\abs{M_{\tau\wedge n}}\le\abs{M_0}+c\tau$, which is integrable, and dominated convergence
applies.
\end{proof}

\begin{nuance}{the conditions cannot be dropped}
Let $S_n$ be the symmetric simple random walk from $0$ and $\tau=\min\{n:S_n=1\}$. Then
$\tau<\infty$ a.s.\ (recurrence), yet $S_\tau=1$ while $\Exp[S_0]=0$. The theorem fails
because $\Exp[\tau]=\infty$ and $S_{\tau\wedge n}$ is unbounded below. The doubling
strategy ``bet $1$, double after every loss, stop at the first win'' always ends with a
profit of $1$, but only by risking unbounded losses along the way. Optional stopping
formalises why there is no free lunch under a finite credit limit or a bounded time.
\end{nuance}

\subsection*{Applications}

\begin{example}[Gambler's ruin by martingales]\label{ex:ruinmart}
For the symmetric walk from $k$ with $\tau=\min\{n:S_n\in\{0,N\}\}$:
$\tau<\infty$ a.s.\ and $\Exp[\tau]<\infty$, because from any state the walk exits within
$N$ steps with probability at least $2^{-N}$, so $\Prob(\tau>mN)\le(1-2^{-N})^m$ and
$\Exp[\tau]=\sum_n\Prob(\tau>n)<\infty$. The stopped walk is bounded by $N$, so
Theorem~\ref{thm:ost}(b) gives $k=\Exp[S_\tau]=N\cdot\Prob(S_\tau=N)$, i.e.\ $h_k=k/N$. For the
duration, $(S_n-k)^2-n$ is a martingale (Example~\ref{ex:martingales}(2) applied to the
increments), so $\Exp[(S_{\tau\wedge n}-k)^2]=\Exp[\tau\wedge n]$ for every $n$. Let
$n\to\infty$: $\tau\wedge n\uparrow\tau$ (monotone convergence on the right) and
$S_{\tau\wedge n}$ is bounded by $N$ (bounded convergence on the left), so
\[
\Exp[\tau]=\Exp[(S_\tau-k)^2]=(N-k)^2\frac kN+k^2\Bigl(1-\frac kN\Bigr)=k(N-k),
\]
confirming Theorem~\ref{thm:duration}. For the biased walk use the martingale
$(q/p)^{S_n}$: with $r=q/p$, $r^k=h\,r^N+(1-h)$, so $h=(1-r^k)/(1-r^N)$, again agreeing with
Theorem~\ref{thm:ruin}.
\end{example}

\begin{theorem}[Wald's identity]\label{thm:wald}
Let $X_1,X_2,\dots$ be i.i.d.\ with $\Exp\abs{X_1}<\infty$ and mean $\mu$, and let $\tau$ be a
stopping time with $\Exp[\tau]<\infty$. Then
\[
\Exp\Bigl[\sum_{i=1}^\tau X_i\Bigr]=\mu\,\Exp[\tau].
\]
\end{theorem}
\begin{proof}
By optional stopping for the bounded stopping time $\tau\wedge n$ applied to the martingale
$S_k-k\mu$: $\Exp[S_{\tau\wedge n}]=\mu\Exp[\tau\wedge n]$. As $n\to\infty$, the right side
tends to $\mu\Exp[\tau]$ by monotone convergence. For the left side, $S_{\tau\wedge n}\to S_\tau$
and $\abs{S_{\tau\wedge n}}\le\sum_{k=1}^\tau\abs{X_k}$, whose expectation is
\[
\sum_{k\ge1}\Exp\bigl[\abs{X_k}\indic_{\{\tau\ge k\}}\bigr]=\sum_{k\ge1}\Exp\abs{X_1}\Prob(\tau\ge k)
=\Exp\abs{X_1}\Exp[\tau]<\infty,
\]
since $\{\tau\ge k\}$ is determined by $X_1,\dots,X_{k-1}$ and so is independent of $X_k$.
Dominated convergence finishes the proof.
\end{proof}

\begin{example}[Expected waiting time for a pattern]\label{ex:pattern}
Toss a fair coin until a given pattern of length $m$ appears; let $\tau$ be the first time it
does. \emph{Claim:}
\[
\Exp[\tau]=\sum_{k=1}^m2^k\,\indic\{\text{the first }k\text{ letters of the pattern equal its last }k\text{ letters}\}.
\]
\emph{Proof.} Before each toss $j$ a new gambler arrives with a fortune of $1$ rupee and bets
it all on toss $j$ showing the first letter of the pattern. If he loses he is out; if he wins
his fortune doubles to $2$ and he bets all of it on toss $j+1$ showing the second letter, and
so on until the pattern is completed (every bet is fair). Let $M_n$ be the total fortune of
all gamblers present up to time $n$ minus the total amount put in, $n$: it is a martingale
(a sum of fair bets) with increments bounded by $2^m$, and $\Exp[\tau]<\infty$ (the pattern
appears in each block of $m$ tosses with probability $2^{-m}$). By
Theorem~\ref{thm:ost}(c), $\Exp[M_\tau]=0$, i.e.\ $\Exp[\text{total fortune at }\tau]=\Exp[\tau]$.

At time $\tau$ the pattern has just been completed, so the gambler who started $m$ tosses
earlier holds $2^m$. A gambler who started $k<m$ tosses earlier is still alive exactly when
the last $k$ tosses equal the first $k$ letters of the pattern, in which case he holds $2^k$;
all others hold $0$. The last $k$ tosses are the last $k$ letters of the pattern, so the
condition is the overlap condition in the claim, and the total fortune is the stated sum.
\hfill$\square$

Examples: HH $\to2^1+2^2=6$ (overlaps at $k=1$ and $k=2$); HT $\to2^2=4$ (only the
full pattern overlaps; this matches Example~\ref{ex:hhht}); HHH $\to2+4+8=14$;
HTH $\to2+8=10$; TTH $\to8$. For HHTHH the overlaps are at $k=1$ (H $=$ H), $k=2$
(HH $=$ HH) and $k=5$, giving $2+4+32=38$.
\end{example}

\section{Maximal and concentration inequalities}

\begin{theorem}[Doob's maximal inequality]\label{thm:doobmax}
If $M$ is a non-negative submartingale, then for $\lambda>0$,
$\Prob\bigl(\max_{k\le n}M_k\ge\lambda\bigr)\le\Exp[M_n]/\lambda$.
\end{theorem}
\begin{proof}
Let $A_k=\{M_k\ge\lambda,M_j<\lambda\ \text{for }j<k\}$, which is determined by $X_1,\dots,X_k$;
the $A_k$ are disjoint with union $A=\{\max_{k\le n}M_k\ge\lambda\}$. By the submartingale
property $\Exp[M_n\given\F_k]\ge M_k$, so
$\Exp[M_n\indic_{A_k}]=\Exp[\indic_{A_k}\Exp[M_n\given\F_k]]\ge\Exp[M_k\indic_{A_k}]\ge\lambda\Prob(A_k)$.
Summing over $k$: $\Exp[M_n]\ge\Exp[M_n\indic_A]\ge\lambda\Prob(A)$.
\end{proof}

(Kolmogorov's inequality, Theorem~\ref{thm:kolmax}, is the special case $M_k=S_k^2$, which is
a non-negative submartingale for mean-zero independent increments.)

\begin{theorem}[Azuma--Hoeffding inequality]\label{thm:azuma}
Let $M$ be a martingale with $\abs{M_k-M_{k-1}}\le c_k$ for constants $c_k$. Then for
$t>0$,
\[
\Prob(M_n-M_0\ge t)\le\exp\Bigl(-\frac{t^2}{2\sum_{k=1}^nc_k^2}\Bigr),
\]
and the same bound holds for $\Prob(M_n-M_0\le-t)$.
\end{theorem}
\begin{proof}
Let $D_k=M_k-M_{k-1}$. Given $\F_{k-1}$, $D_k$ has conditional mean $0$ and lies in
$[-c_k,c_k]$, an interval of length $2c_k$. Hoeffding's lemma
(Lemma~\ref{lem:hoeffding}), applied to the conditional distribution, gives
$\Exp[e^{sD_k}\given\F_{k-1}]\le e^{s^2(2c_k)^2/8}=e^{s^2c_k^2/2}$. Iterating the tower
property,
\[
\Exp[e^{s(M_n-M_0)}]=\Exp\bigl[e^{s(M_{n-1}-M_0)}\Exp[e^{sD_n}\given\F_{n-1}]\bigr]
\le e^{s^2c_n^2/2}\Exp[e^{s(M_{n-1}-M_0)}]\le\dots\le e^{s^2\sum c_k^2/2}.
\]
By Markov's inequality, $\Prob(M_n-M_0\ge t)\le e^{-st+s^2\sum c_k^2/2}$; the best $s$ is
$t/\sum c_k^2$, giving the bound. The lower tail follows by applying this to $-M$.
\end{proof}

\begin{example}[Bounded differences]
Let $f(x_1,\dots,x_n)$ change by at most $c_k$ when only the $k$th coordinate changes, and let
$X_i$ be independent. The Doob martingale $M_k=\Exp[f(X)\given X_1,\dots,X_k]$ has
$\abs{M_k-M_{k-1}}\le c_k$ (resampling $X_k$ changes $f$ by at most $c_k$ in conditional
expectation), $M_0=\Exp[f(X)]$ and $M_n=f(X)$. Azuma gives McDiarmid's inequality:
$\Prob(f(X)-\Exp f(X)\ge t)\le e^{-t^2/(2\sum c_k^2)}$ (a sharper constant, $e^{-2t^2/\sum c_k^2}$,
holds with a more careful argument). Application: the number $D$ of distinct values among $n$
uniform draws from $m$ boxes changes by at most $1$ per draw, so
$\Prob(\abs{D-\Exp D}\ge t)\le2e^{-t^2/(2n)}$.
\end{example}

\section{Convergence of martingales}

\begin{theorem}[Doob's upcrossing inequality]\label{thm:upcross}
Let $M$ be a martingale, $a<b$, and let $U_n$ be the number of upcrossings of $[a,b]$ by
$M_0,\dots,M_n$ (passages from below $a$ to above $b$). Then
$(b-a)\Exp[U_n]\le\Exp[(M_n-a)^-]$, where $x^-=\max(-x,0)$.
\end{theorem}
\begin{proof}
Define a predictable betting strategy: $H_k=1$ (``hold one unit'') from the first time after
$M$ drops below $a$ until the next time it exceeds $b$, and $H_k=0$ otherwise. Formally
$H_1=\indic_{\{M_0<a\}}$ and $H_k=\indic_{\{H_{k-1}=1,M_{k-1}\le b\}}+\indic_{\{H_{k-1}=0,M_{k-1}<a\}}$;
$H_k$ depends on $M_0,\dots,M_{k-1}$, so it is predictable. Each completed upcrossing earns at
least $b-a$ in $(H\cdot M)$. A final incomplete stretch (bought below $a$ and not yet sold
above $b$) can lose at most $(M_n-a)^-$. Hence
$(H\cdot M)_n\ge(b-a)U_n-(M_n-a)^-$. By Theorem~\ref{thm:transform},
$\Exp[(H\cdot M)_n]=0$, and the claim follows.
\end{proof}

\begin{theorem}[Martingale convergence theorem]\label{thm:martconv}
If $M$ is a martingale with $\sup_n\Exp\abs{M_n}<\infty$, then $M_n$ converges almost surely
to a finite limit $M_\infty$. In particular every non-negative martingale converges almost
surely.
\end{theorem}
\begin{proof}
Let $C=\sup_n\Exp\abs{M_n}$. By the upcrossing inequality,
$\Exp[U_n]\le(C+\abs a)/(b-a)$ for all $n$, so the total number of upcrossings
$U_\infty=\lim U_n$ has finite expectation, hence is finite almost surely. A real sequence
that has only finitely many upcrossings of every rational interval $[a,b]$ cannot oscillate
(if $\liminf<\limsup$ there would be rationals $a<b$ between them with infinitely many
upcrossings), so $M_n$ converges in $[-\infty,\infty]$ almost surely. By Fatou's lemma
(Chapter~\ref{ch:measure}), $\Exp\abs{M_\infty}\le\liminf\Exp\abs{M_n}\le C<\infty$, so the limit
is finite almost surely. For $M\ge0$, $\Exp\abs{M_n}=\Exp[M_n]=\Exp[M_0]$ is constant.
\end{proof}

\begin{example}[P\'olya's urn limit]\label{ex:polya}
Start the urn with one red and one blue ball. After $n$ draws, let $K$ be the number of red
balls added. Any particular sequence with $k$ red draws has probability
$\frac{k!\,(n-k)!}{2\cdot3\cdots(n+1)}=\frac{k!(n-k)!}{(n+1)!}$ (the numerators $1,2,\dots,k$ for
red and $1,\dots,n-k$ for blue, denominators $2,3,\dots,n+1$), independent of the order. Since
there are $\binom nk$ such sequences, $\Prob(K=k)=\binom nk\frac{k!(n-k)!}{(n+1)!}=\frac1{n+1}$:
\emph{uniform on $\{0,\dots,n\}$}. The fraction red is $F_n=(K+1)/(n+2)$, which converges in
distribution to $U(0,1)$. Since $F_n$ is a bounded martingale it converges almost surely by
Theorem~\ref{thm:martconv}; the limit therefore has the uniform law. So each run of the urn
settles to a random proportion, uniformly distributed, and that proportion is determined by
the early draws.
\end{example}

\begin{nuance}{convergence a.s. need not mean convergence in $L^1$}
The martingale $M_n=\prod_{i\le n}Y_i$ with $Y_i$ i.i.d., $\Prob(Y=0)=\Prob(Y=2)=\tfrac12$,
has $\Exp[M_n]=1$ for all $n$, but $M_n\to0$ almost surely (the first zero factor kills it).
So $\Exp[M_\infty]=0\ne1$. Almost sure convergence of an $L^1$-bounded martingale does not
imply convergence of expectations; this requires uniform integrability. This is the same
phenomenon as in Example~\ref{ex:conversefail}(ii).
\end{nuance}

\section{Optimal stopping}

Suppose rewards $g_0,g_1,\dots,g_N$ are revealed over time (adapted), and we may choose a
stopping time $\tau\le N$ to collect $g_\tau$. What is the best expected reward, and what
is the best stopping rule?

\begin{theorem}[Dynamic programming for optimal stopping]\label{thm:snell}
Define $V_N=g_N$ and $V_n=\max\bigl(g_n,\Exp[V_{n+1}\given\F_n]\bigr)$ for $n<N$ (the value
of the best strategy from time $n$ on). Then $\sup_{\tau\le N}\Exp[g_\tau]=\Exp[V_0]$, and the
supremum is attained by $\tau^*=\min\{n:g_n=V_n\}$ (stop as soon as the reward at least
equals the value of continuing).
\end{theorem}
\begin{proof}
By construction $V_n\ge g_n$ and $V_n\ge\Exp[V_{n+1}\given\F_n]$, so $V$ is a supermartingale
dominating $g$. For any stopping time $\tau\le N$, $V_{\tau\wedge n}$ is a supermartingale
(the transform argument of Theorem~\ref{thm:ost0} with inequality), so
$\Exp[g_\tau]\le\Exp[V_\tau]\le\Exp[V_0]$. For $\tau^*$: before $\tau^*$ we have
$V_n>g_n$, so $V_n=\Exp[V_{n+1}\given\F_n]$, which means $V_{\tau^*\wedge n}$ is a martingale
and $\Exp[V_{\tau^*}]=\Exp[V_0]$; also $V_{\tau^*}=g_{\tau^*}$ (at time $N$, $V_N=g_N$).
Hence $\Exp[g_{\tau^*}]=\Exp[V_0]$.
\end{proof}

\begin{example}[Rolling a die up to three times]\label{ex:diceopt}
You roll a fair die up to $3$ times and may stop at any roll, taking its value as your
reward. Work backwards. With one roll left the value is $V_1=\Exp[X]=3.5$. With two rolls
left you stop on the first roll if it exceeds $3.5$:
$V_2=\Exp[\max(X,3.5)]=\frac{3\cdot3.5+4+5+6}6=4.25$. With three rolls left you stop on the
first roll if it exceeds $4.25$ (i.e.\ is at least $5$):
$V_3=\Exp[\max(X,4.25)]=\frac{4\cdot4.25+5+6}6=\frac{28}6\approx4.667$.
\end{example}

\begin{theorem}[The secretary problem]\label{thm:secretary}
$n$ candidates, all with distinct quality, arrive in uniformly random order. After each
interview you see only the \emph{relative rank} of the candidate among those seen so far, and
must immediately accept or reject; a rejected candidate cannot be recalled. To maximise the
probability of choosing the overall best candidate, use the rule: reject the first $r-1$
candidates, then accept the first candidate who is the best so far. The success probability is
\[
P_n(r)=\frac{r-1}n\sum_{k=r}^n\frac1{k-1}\qquad(r\ge2),
\]
the optimal $r$ is approximately $n/e$, and the optimal success probability tends to
$1/e\approx0.368$ as $n\to\infty$.
\end{theorem}
\begin{proof}
\emph{Success probability.} The rule picks the overall best candidate exactly when, for some
position $k\ge r$, the best candidate is at position $k$ and the best of the first $k-1$
candidates is among the first $r-1$ (so that no one between positions $r$ and $k-1$ was
the best so far). The best candidate is at position $k$ with probability $1/n$, and
independently the best among the first $k-1$ is in the first $r-1$ positions with probability
$(r-1)/(k-1)$ (by exchangeability of the relative order). Summing over $k$ gives $P_n(r)$.

\emph{Optimality of threshold rules.} Stopping at a best-so-far candidate at time $k$ wins with
probability $k/n$ (the best of the first $k$ is the overall best). The value $c_k$ of
continuing optimally after time $k$ is non-increasing in $k$ (fewer chances remain), while
$k/n$ is increasing. So the set of $k$ at which stopping (given a best-so-far) is optimal is
$\{k:k/n\ge c_k\}$, an upper interval $\{r,r+1,\dots,n\}$, by
Theorem~\ref{thm:snell} applied with state ``this candidate is best so far''. The optimal rule
is therefore a threshold rule of the form above, and it remains to optimise over $r$.

\emph{Asymptotics.} Put $x=r/n$. Then
$P_n(r)=\frac{r-1}n\sum_{k=r}^n\frac1{k-1}\to x\int_x^1\frac{du}u=-x\ln x$. This is maximised
where $\frac d{dx}(-x\ln x)=-\ln x-1=0$, i.e.\ $x=1/e$, with maximum value $1/e$.
\end{proof}

\begin{interview}{the secretary rule and its cousins}
``Look at the first $37\%$, then leap at the first one better than everything seen'' wins
with probability about $37\%$, remarkably independent of $n$. Variants: maximise the expected
rank (the answer is a different threshold sequence), allow a probability of recall, or let
the number of candidates be unknown. In every variant, work backwards from the last stage:
$V_n=\max(\text{stop},\text{continue})$.
\end{interview}

\section*{Exercises}
\addcontentsline{toc}{section}{Exercises}

\begin{enumerate}
\item Show that for a symmetric simple random walk $S_n^3-3nS_n$ is a martingale and
use it with optional stopping to compute $\Exp[\tau]$ and
$\Exp[\tau\mid S_\tau=N]$ for the walk started at $0$ and stopped on hitting $\pm N$.
\item Show that $\Exp[\tau]=N^2$ for the symmetric walk from $0$ stopped on exiting
$(-N,N)$.
\item Use Example~\ref{ex:pattern} to find the expected number of tosses to get HHTHH,
THT and HTHT.
\item Two players keep tossing a fair coin; $A$ wins when HH first appears, $B$ wins when
HT first appears. Show that $B$ wins with probability $\tfrac12$ (in fact, they alternate
on the same sequence). Then for the patterns HH vs.\ TH show who is more likely to win.
\item Prove that if $M$ is a martingale and $\varphi$ is convex with $\Exp\abs{\varphi(M_n)}<\infty$
then $\varphi(M_n)$ is a submartingale. (Jensen for conditional expectation.)
\item A fair die is rolled up to $4$ times with the option to stop. Find the value of the
optimal stopping problem and the optimal rule.
\item In the secretary problem with $n=10$, compute $P_{10}(r)$ for each $r$ and find the best
threshold.
\item (Challenge) Show that the condition $\Exp[\tau]<\infty$ in
Theorem~\ref{thm:ost}(c) cannot be dropped: give a martingale with bounded increments and
an a.s.\ finite stopping time $\tau$ with $\Exp[\tau]=\infty$ and $\Exp[M_\tau]\ne\Exp[M_0]$.
\end{enumerate}
